\documentclass{article}
\begin{document}
\section*{Compatibility with changing payload presheaf}
An arbitrary ordinary natural transformation of category-valued payload presheaves acts in two independently defined ways. On the iterate it induces a native Grothendieck functor at each stage. On a decorated chain square it fixes the grade square, maps its terminal payload arrow, and uses naturality on the actual terminal base arrow to identify the required target. Both sides preserve identities and binary composition. Induction compares the entire realized arrow, retaining its head and recursive tail. Faithfulness and object injectivity of the fibre comparison then prove that source payload change commutes with source base restriction. The length-indexed coproduct inherits both squares; both payload-variable constructions are ordinary functors with the specified output universes. Their pointwise categorical equivalences have forward components strictly natural in both variables. No pseudonatural input, common-universe bundled natural isomorphism, span-equipment sequential product identification, flattening, or free-monad theorem claimed.
\par Manuscript: equations (groth) and (telescope), pinned commit 458ff36a4798e432dee5947afc402f9903a47f14. Baseline: merged 180. Fresh project-source Lean 4.35 replay; external dependency caches reused. Independent scoped review passed; see the separate review and replay logs.
\end{document}
